\documentclass[10pt,xcolor=pst,aspectratio=169]{beamer}

\usepackage{etex}

%\usetheme{Boadilla}
%\usecolortheme{wolverine}
\usecolortheme{dolphin}
%\setbeamercovered{transparent}
%\setbeamercolor{block body}{bg=yellow}

\addtobeamertemplate{navigation symbols}{}{%
\usebeamerfont{footline}%
\usebeamercolor[fg]{footline}%
\hspace{1em}%
\insertframenumber/\inserttotalframenumber
}

\usepackage[utf8]{inputenc}
\usepackage[english,russian]{babel}
\usepackage[OT1]{fontenc}
\usepackage{amsmath}
\usepackage{amsfonts}
\usepackage{amssymb}
\usepackage{graphicx}
\usepackage{wrapfig}
\usepackage[3D]{movie15}
\usepackage{animate}
\usepackage{ragged2e}
\usepackage{listings}
\usepackage{color}
\usepackage{pst-all}

\usepackage{tikz}
\usetikzlibrary{
	mindmap,
	arrows, % стрелки
	shapes.misc, % фигуры
	chains, % цепочки
	positioning, % позиционирование элементов
	scopes, % создание дополнительных веток
	shadows % тени
	}

\graphicspath{{pic/}}

\author{\textbf{Губкин А.С.}}

\title[Численные методы в физике]{Численные методы в физике}

\logo{\includegraphics[width=0.1\linewidth]{LOGO_2.EPS}}

\institute[ТюмФ ИТПМ СО РАН]{Тюменский филиал Института теоретической и прикладной механики\\ им. С. А. Христиановича СО РАН, г. Тюмень}

%\date{6 октября 2015 г.}

\begin{document}

\lstset{ %
	language=[ANSI]C++,                 % выбор языка для подсветки (здесь это С++)
	keywordstyle=\color{blue},
	commentstyle=\color{gray},
	basicstyle=\scriptsize,
%basicstyle=\small\sffamily, % размер и начертание шрифта для подсветки кода
	numbers=left,               % где поставить нумерацию строк (слева\справа)
	numberstyle=\tiny,           % размер шрифта для номеров строк
%stepnumber=1,                   % размер шага между двумя номерами строк
	numbersep=4pt,                % как далеко отстоят номера строк от подсвечиваемого кода
%backgroundcolor=\color{white}, % цвет фона подсветки - используем \usepackage{color}
	showspaces=false,            % показывать или нет пробелы специальными отступами
	showstringspaces=false,      % показывать или нет пробелы в строках
	showtabs=false,             % показывать или нет табуляцию в строках
	frame=single,              % рисовать рамку вокруг кода
%tabsize=2,                 % размер табуляции по умолчанию равен 2 пробелам
	captionpos=t,              % позиция заголовка вверху [t] или внизу [b] 
	breaklines=true,           % автоматически переносить строки (да\нет)
	breakatwhitespace=true, % переносить строки только если есть пробел
	escapeinside={\%*}{*)}   % если нужно добавить комментарии в коде
}

%SLIDE #
\begin{frame}

	\transdissolve[duration=0.1]
	\titlepage

\end{frame}

%SLIDE #
\begin{frame}{Диссипация и дисперсия численного решения}

	\transdissolve[duration=0.1]
	\justifying
	\large

	Определим диссипацию и диспесиию для дифференциального волнового уравнения. Возьмем решение в виде:

	\[
		u(x, t) = u_{0} \exp{(i (\omega t - k x))},
	\]

	где $\omega = 2 \pi \nu$ -- круговая частота, $k = 2 \pi / \lambda$ -- волновое число.\\

	Подставив это решение в волновое уравнение, получим зависимость $\omega = \omega(k)$, которая называется \textbf{дисперсионным соотношением}.\\

	Если $\omega$ -- комплексное число, волна затухает!

	\[
		\exp{((- Im \omega) t)} = \exp{(- \gamma t)}.
	\]

\end{frame}

%SLIDE #
\begin{frame}{Фазовая и групповая скорость}

	\transdissolve[duration=0.1]
	\justifying
	\large

	\textbf{Фазовая скорость} -- это скорость, с которой движется фаза или отдельной гармоники:

	\[
		\frac{Re \omega}{k} = c_{f}.
	\]

	\textbf{Групповая скорость} -- это скорость волнового пакета, состоящего из гармонических волн с близкими волновыми числами (Передача энергии осуществляется с групповой скоростью!):

	\[
		\frac{d}{d k} Re \omega = c_{g}.
	\]

\end{frame}

%SLIDE #
\begin{frame}{Дисперсия волн}

	\transdissolve[duration=0.1]
	\justifying
	\large

	Если фазовая/групповая скорость зависит от $k$, то гармоники с разными волновыми числами распространяются с разными скоростями. Такое явление называется \textbf{дисперсией}.

\end{frame}

%SLIDE #
\begin{frame}{Пример №1}

	\transdissolve[duration=0.1]
	\justifying
	\large

	\[
		\frac{\partial^{2} u}{\partial t^{2}} = c^{2} \frac{\partial^{2} u}{\partial x^{2}} \Rightarrow - \omega^{2} = - c^{2} k^{2} \Rightarrow \omega = c k \Rightarrow \frac{Re \omega}{k} = \frac{d}{d k} Re \omega = c.
	\]
	
	Точное решение -- волна без дисперсии и затухания.

\end{frame}

%SLIDE #
\begin{frame}{Пример №2}

	\transdissolve[duration=0.1]
	\justifying
	\large

	\[
		\frac{\partial u}{\partial t} + c \frac{\partial u}{\partial x} = \mu \frac{\partial^{2} u}{\partial x^{2}} \Rightarrow i \omega - i c k = - \mu k^{2} \Rightarrow \omega(k) = c k + i \mu k^{2}.
	\]
	
	Точное решение -- затухающая недиспергирующая волна.

\end{frame}

\end{document}